\item \textbf{Stating the hypotheses.} In each of the following situations a significance test for a population mean $\mu$ is called for.  Identify the null hypothesis $H_0$ and the alternative hypothesis $H_A$ in each case. Make sure to use proper notation! (\textit{e.g. $H_0: \mu = \mu_0$}).
	\begin{enumerate}
		\item 50 college students are selected to participate in a Kinesiology study. A researcher believes his workout regime increases body mass. All 50 participate in a weight lifting regime three times a week and their before and after body masses are measured. 

	\item Ten years ago there was a Mercury spill in Lake Tomahawk that caused serious damage to humans and wildlife. Today, researchers want to know whether it is safe for humans to use the lake recreationally. The EPA says water is safe if no more than 2 mg/L mercury is in the water. The researchers calculated an average of 1.8 mg/L based on multiple measurements from the lake.

	\item A public accounting firm is auditing a very large corporation. One of the components the auditors need to check is the balance on the company's accounts receivable at fiscal year end. It would be very costly to reconcile every customer of the company and how much they owe. A partner at the accounting firm claims the average accounts receivable balance is \$159,000. A junior associate at the accounting firm believes the partner is wrong, but would like to prove with statistical evidence. Based on a sample, the junior associate calculated the mean accounts receivable balance to be \$164,000.

		\item The population mean IQ is 105. A psychologist wants to test the hypothesis that a new cognitive treatment increases the IQ of patients who undergo the treatment. The pyschologist trys the new treatment on 40 people and calculates an average IQ after the treatment of 106.7.
	\end{enumerate}